\chapter{INTRODUCTION}

\section{Background}

The rapid development of high-throughput genomic technologies has made it possible to measure the expression levels of thousands of genes simultaneously. Gene expression data can therefore provide useful information for distinguishing between different biological and disease states. Machine learning classification techniques can be applied to such data to identify patterns that may help distinguish between classes of samples.

A major challenge in the analysis of gene expression data is its high dimensionality. The number of measured features can be considerably larger than the number of available samples, resulting in a small-sample, high-dimensional setting. In this setting, the feature space may contain irrelevant or redundant variables, while the limited number of observations can make reliable estimation of classification models difficult. These characteristics can increase computational requirements and may contribute to overfitting and poor generalization.

Feature selection methods can be used in 
data pre-processing to achieve efficient data reduction\citep{8677302}
Different feature selection methods use different criteria to determine which features should be retained. These differences can result in feature subsets of different sizes and can influence the subsequent classification performance.

In high-dimensional gene expression classification, it is therefore important to consider not only whether the number of features is reduced, but also how the selected features affect classification performance. A reduction in dimensionality does not automatically imply an improvement in predictive performance. The effect may depend on characteristics of the dataset, the feature selection method, the distribution of the classes, and the classification procedure used.

This study focuses on high-dimensional human gene expression microarray data and investigates the effect of different feature selection techniques on classification performance. Two publicly available human gene expression datasets with different characteristics are considered. The study compares feature selection approaches from different methodological categories and examines both the resulting feature subsets and their classification performance.


\section{Problem Statement}

High-dimensional gene expression datasets contain a large number of measured features relative to the number of available samples. This ``small $n$, large $p$'' structure presents challenges for classification because many of the available features may be irrelevant or redundant, while the limited number of observations makes it difficult to estimate complex relationships reliably. As a result, classification models may experience overfitting, increased computational cost, and reduced generalization when trained using the complete feature space.

Feature selection provides a possible approach for addressing these challenges by identifying a smaller subset of informative features. However, different feature-selection techniques use different criteria for determining which features are retained. Consequently, they may produce subsets of substantially different sizes and may have different effects on classification performance. Feature selection may improve, maintain, or reduce predictive performance depending on the characteristics of the dataset and the method employed.

This creates a need to empirically investigate the effect of different feature-selection techniques on high-dimensional gene expression classification. In particular, it is necessary to determine whether feature selection consistently improves classification performance across datasets, whether different selection methods produce different outcomes, and how the dimensionality of the resulting feature subsets relates to classification performance.

Gene expression microarray datasets contain thousands of probes but relatively few samples. This high dimensionality can negatively affect classification performance by increasing model variance, promoting overfitting, introducing noise from irrelevant or redundant features, and reducing generalization in classification tasks. While feature selection is commonly used to address this issue, it is not always clear which feature selection techniques most effectively improve classification performance on high-dimensional gene expression data, nor whether their effect is consistent across datasets with different characteristics such as class balance.


\section{Aim and Objectives}

\subsection{Aim}


The aim of this study is to investigate the effect of different feature-selection techniques on the classification performance of high-dimensional human gene expression microarray data.

\subsection{Objectives}

The specific objectives of the study are to:

\begin{enumerate}
\item To examine the structure and characteristics of high-dimensional human gene expression datasets used for classification.

\item To implement selected filter, wrapper, and model-based feature-selection techniques for identifying informative features.

\item To train machine learning classification models using features selected by each method.

\item To investigate the relationship between the number of selected features and classification performance across the datasets.

\item To assess whether the effect of feature selection is consistent across datasets with different characteristics.


\end{enumerate}

\section{Significance of the Study}
Understanding the effect of feature selection is important when applying machine learning methods to high-dimensional gene expression data. Selecting a smaller set of features can reduce the computational burden associated with training classification models and can make the resulting models easier to analyze. However, inappropriate feature selection may also remove useful information and negatively affect predictive performance.

This study provides a comparative evaluation of different feature selection approaches in a common classification framework. By considering datasets with different characteristics, the study examines whether the effect of feature selection is consistent across classification problems rather than assuming that a particular method will always produce better performance.

The study also considers the relationship between the number of features selected and classification performance. This provides a mathematical and computational perspective on feature selection and contributes to understanding how changes in the feature space can influence classification outcomes in high-dimensional gene expression data.

\section{Scope of the Study}
The study is restricted to human gene expression microarray data obtained from the GEO. These datasets contain gene expression measurements with a large number of genes relative to the number of samples.

The study focuses on the application and comparison of selected feature selection techniques. The selected features are subsequently used for classification, and the resulting performance is compared with classification using the full feature space.

The study focuses on classification performance and the dimensionality of the resulting feature subsets. It does not attempt to establish biological significance for individual selected probes or provide a biological interpretation of the genes beyond what is necessary to describe the datasets.

\section{Organisation of the Thesis}
This thesis is organized into six chapters. 

\begin{itemize}
    \item Chapter 1 introduces the study by presenting the background, problem statement, significance, scope, and aims and objectives of the research.

\item Chapter 2 presents a review of the literature related to high-dimensional gene expression data, feature selection, the major categories of feature selection methods, and their application to classification problems. It also identifies the research gap addressed by the study.

\item Chapter 3 describes the datasets, data preparation procedures, feature selection methods, classification procedure, and evaluation framework used in the study.

\item Chapter 4 presents the results obtained from the analysis of the datasets, including the characteristics of the selected feature subsets.

\item Chapter 5 discusses the classification results and examines the effect of feature selection on classification performance across the datasets.

\item Chapter 6 presents the conclusions drawn from the study and provides recommendations for future work.
\end{itemize}
